\documentclass[a4paper,12pt]{extarticle}
\usepackage{geometry}
\usepackage[T1]{fontenc}
\usepackage[utf8]{inputenc}
\usepackage[english,russian]{babel}
\usepackage{amsmath}
\usepackage{amsthm}
\usepackage{amssymb}
\usepackage{fancyhdr}
\usepackage{setspace}
\usepackage{graphicx}
\usepackage{colortbl}
\usepackage{tikz}
\usepackage{pgf}
\usepackage{subcaption}
\usepackage{listings}
\usepackage{indentfirst}
\usepackage[
backend=biber,
style=numeric,
maxbibnames=99
]{biblatex}
\addbibresource{refs.bib}
\usepackage[colorlinks,citecolor=blue,linkcolor=blue,bookmarks=false,hypertexnames=true, urlcolor=blue]{hyperref} 
\usepackage{indentfirst}
\usepackage{mathtools}
\usepackage{booktabs}
\usepackage[flushleft]{threeparttable}
\usepackage{tablefootnote}

\usepackage{chngcntr} % нумерация графиков и таблиц по секциям
\counterwithin{table}{section}
\counterwithin{figure}{section}

\graphicspath{{graphics/}}%путь к рисункам

\makeatletter
% \renewcommand{\@biblabel}[1]{#1.} % Заменяем библиографию с квадратных скобок на точку:
\makeatother

\geometry{left=2.5cm}% левое поле
\geometry{right=1.0cm}% правое поле
\geometry{top=2.0cm}% верхнее поле
\geometry{bottom=2.0cm}% нижнее поле
\setlength{\parindent}{1.25cm}
\renewcommand{\baselinestretch}{1.5} % междустрочный интервал


\newcommand{\bibref}[3]{\hyperlink{#1}{#2 (#3)}} % biblabel, authors, year
\addto\captionsrussian{\def\refname{Список литературы (или источников)}} 

\renewcommand{\theenumi}{\arabic{enumi}}% Меняем везде перечисления на цифра.цифра
\renewcommand{\labelenumi}{\arabic{enumi}}% Меняем везде перечисления на цифра.цифра
\renewcommand{\theenumii}{.\arabic{enumii}}% Меняем везде перечисления на цифра.цифра
\renewcommand{\labelenumii}{\arabic{enumi}.\arabic{enumii}.}% Меняем везде перечисления на цифра.цифра
\renewcommand{\theenumiii}{.\arabic{enumiii}}% Меняем везде перечисления на цифра.цифра
\renewcommand{\labelenumiii}{\arabic{enumi}.\arabic{enumii}.\arabic{enumiii}.}% Меняем везде перечисления на цифра.цифра

\begin{document}
\input{title_kr}% это титульный лист - выберите подходящий вам из имеющихся в проекте вариантов (kr - курсовая работа у 3 курса, vkr - выпускная квалификационная работа у 4 курса)
\newpage
\setcounter{page}{2}

{
	\hypersetup{linkcolor=black}
	\tableofcontents
}

\newpage

\newpage
\section*{Аннотация}   % this is how to use russian
Работа посвящена разработке программной платформы для исследования методов решения задачи коммивояжёра (TSP, \emph{Traveling Salesman Problem}) и задачи множественного коммивояжёра в постановке min--max с несколькими депо (Min--Max Multi-Depot mTSP). Платформа поддерживает загрузку наборов данных (графов), запуск набора известных алгоритмов и эвристик, а также их сравнение по качеству решения и времени работы.

В рамках проекта реализуются процедуры построения начальных решений и улучшения на основе локального поиска и гибридных модификаций; планируется предложить новую эвристику или алгоритм, превосходящий выбранные аналоги на тестовых экземплярах. Дополнительно реализуется \emph{anytime}-солвер, который быстро строит допустимое решение и продолжает улучшать его при наличии оставшегося времени.

Отдельно планируется исследовать применимость графовых нейронных сетей (GNN, \emph{Graph Neural Networks}) для поддержки эвристик, например для отбора перспективных рёбер/кандидатов и настройки параметров алгоритмов.

\addcontentsline{toc}{section}{Аннотация}

\section*{Ключевые слова}
Задача коммивояжёра (TSP), Min--Max Multi-Depot mTSP, маршрутизация (routing), эвристики (heuristics), локальный поиск (local search), гибридные методы (hybrid methods), графовые нейронные сети (GNN)
\pagebreak


\section{Введение}

\subsection{Мотивация и цель работы}
Работа направлена на построение единой исследовательской платформы для задач TSP и Min--Max Multi-Depot mTSP: от формализации и реализации алгоритмов до воспроизводимых экспериментов.

\subsection{Вклад работы}
\begin{enumerate}
    \item Реализовано единое C++-ядро для двух постановок маршрутизации.
    \item Реализован набор точных, эвристических и метаэвристических методов.
    \item Реализована Python-обвязка для серийных запусков, валидации и сбора метрик.
\end{enumerate}

\subsection{Структура отчета}
В главах 1--6 последовательно рассмотрены постановки задач, реализованные алгоритмы, архитектура платформы, экспериментальный протокол и результаты, а также опциональное ML/GNN-направление.

\pagebreak

\section{Основные определения, постановки и критерии качества}

\subsection{Основные обозначения и входные данные}
Используются взвешенные графы $G=(V,E)$, где $V$ --- множество вершин, а $w(i,j)$ --- стоимость перехода.
Поддерживаются три формата входа: готовая матрица расстояний, координаты с метрическими функциями, а также формат $\texttt{latlon}$.

\subsection{Формальная постановка задачи TSP}
Для TSP требуется найти гамильтонов цикл минимальной длины:
\begin{equation}
\label{eq:tsp_obj}
\min_{\pi \in S_n}\; \sum_{k=1}^{n} w\bigl(\pi_k,\pi_{k+1}\bigr), \qquad \pi_{n+1}=\pi_1.
\end{equation}

\subsection{Формальная постановка Min--Max Multi-Depot mTSP}
Пусть задано множество депо $D=\{d_1,\dots,d_p\}\subseteq V$.
Необходимо построить набор маршрутов, покрывающих всех клиентов $V\setminus D$, минимизируя длину самого длинного маршрута:
\begin{equation}
\label{eq:minmax_mtsp_obj}
\min\; \max_{r} \; L_r,
\end{equation}
где $L_r$ --- длина $r$-го маршрута, начинающегося и заканчивающегося в соответствующем депо.

\subsection{Критерии оценки алгоритмов}
В работе используются следующие метрики:
\begin{itemize}
    \item качество решения: $\texttt{len}$ для TSP и $\texttt{max\_len}$ для Min--Max mTSP;
    \item вычислительное время;
    \item валидность маршрутов;
    \item устойчивость по серии запусков: среднее значение и лучшее значение.
\end{itemize}

\pagebreak

\section{Алгоритмы решения TSP в реализованной платформе}

\subsection{Базовые эвристики: Nearest Neighbor и 2-opt}
\begin{itemize}
    \item \textbf{Nearest Neighbor}: построение начального тура жадным выбором ближайшей непосещенной вершины.
    \item \textbf{2-opt}: локальное улучшение за счет разворота подотрезков маршрута.
\end{itemize}

\subsection{Точный метод: Held--Karp}
Реализована динамика по подмножествам для получения оптимального решения на малых размерах графа.
Для контроля вычислительной сложности вводится ограничение на максимальное число вершин.

\subsection{Метаэвристики: ILS и GLS}
\begin{itemize}
    \item \textbf{ILS}: пертурбации типа double-bridge + локальный поиск 2-opt.
    \item \textbf{GLS}: штрафы на ребра и итеративное улучшение с параметром $\lambda$.
\end{itemize}

\subsection{Интеграция LKH}
Реализована обертка над внешним LKH: генерация TSPLIB-файлов, формирование параметров запуска, чтение оптимизированного тура и возврат результата в общий пайплайн.

\pagebreak

\section{Алгоритмы решения Min--Max Multi-Depot mTSP}

\subsection{Построение начального решения: Random и NN}
\begin{itemize}
    \item \textbf{Random}: случайная перестановка клиентов и распределение по маршрутам.
    \item \textbf{NN}: жадное добавление клиентов в маршрут с минимальной текущей длиной.
\end{itemize}

\subsection{Локальное улучшение маршрутов: 2-opt}
Для каждого маршрута выполняется локальный поиск 2-opt.
Целевая метрика на уровне решения --- максимум длин маршрутов.

\subsection{Муравьиный алгоритм (Ant Colony)}
Реализован ACO-подход с:
\begin{itemize}
    \item candidate lists;
    \item феромонной матрицей и испарением;
    \item параллельной симуляцией нескольких муравьев;
    \item пост-улучшением найденных маршрутов через 2-opt.
\end{itemize}

\subsection{Сценарии композиции солверов}
Платформа поддерживает последовательный запуск нескольких шагов (\texttt{--step}), что позволяет строить гибридные пайплайны улучшения решения.

\pagebreak

\section{Архитектура программной платформы}

\subsection{C++-ядро: интерфейсы и фабрики}
Реализованы абстрактные интерфейсы \texttt{Solver} для TSP и Min--Max mTSP, а также фабрики для регистрации и создания алгоритмов по имени.

\subsection{Модели данных и парсинг экземпляров}
Реализованы singleton-инстансы задач для TSP и Min--Max mTSP с поддержкой нескольких форматов входных данных и метрик расстояний.

\subsection{Python-оркестратор экспериментов}
Python-раннер отвечает за:
\begin{itemize}
    \item подготовку payload для C++-ядра;
    \item выполнение серий запусков;
    \item валидацию маршрутов;
    \item сохранение итоговых JSON-артефактов.
\end{itemize}

\subsection{Логирование и наблюдаемость}
Реализован потокобезопасный логгер с режимами \texttt{info}/\texttt{debug}, периодическими snapshot-записями, историей улучшений и сохранением best/latest маршрутов.

\pagebreak

\section{Экспериментальная часть и анализ результатов}

\subsection{Протокол экспериментов}
\begin{enumerate}
    \item Выбор набора задач для TSP и Min--Max mTSP.
    \item Фиксация конфигураций и параметров солверов.
    \item Серийные запуски с оценкой среднего качества и времени.
\end{enumerate}

\subsection{Результаты по TSP}
В разделе приводятся сравнительные таблицы по алгоритмам \texttt{nearest}, \texttt{2-opt}, \texttt{heldkarp}, \texttt{ils}, \texttt{gls}, \texttt{lkh}.
Рекомендуется показывать: итоговую длину, время, а также вариативность по сериям запусков.

\subsection{Результаты по Min--Max Multi-Depot mTSP}
Приводится сравнение \texttt{random}, \texttt{nn}, \texttt{two-opt}, \texttt{ant} по метрикам $\max L_r$, времени и стабильности.

\subsection{Обсуждение}
Фиксируются практические компромиссы:
\begin{itemize}
    \item качество vs время;
    \item влияние локального поиска;
    \item влияние candidate lists и параллелизма в ACO.
\end{itemize}

\pagebreak

\section{ML/GNN-направление как расширение платформы (опционально)}

\subsection{Текущее состояние}
В репозитории присутствует отдельный блок NeuroLKH с обучающими скриптами и предобученными моделями, который рассматривается как база для дальнейшей интеграции.

\subsection{Потенциальная интеграция}
Перспективные направления:
\begin{itemize}
    \item предсказание candidate edges для ускорения локального поиска;
    \item параметрическая адаптация эвристик;
    \item гибридные схемы ML + классические солверы.
\end{itemize}

\subsection{План валидации ML-гипотез}
Для оценки интеграции предлагаются A/B-сравнения (с ML и без ML) на общем протоколе экспериментов и одинаковых наборах задач.

\newpage 
\printbibliography[heading=bibintoc] 

% % \begin{thebibliography}{0}
% % 	\bibitem{chirkova18}\hypertarget{chirkova18}{}
% % 	\href{https://arxiv.org/abs/1810.10927}
% % 	{Nadezhda Chirkova, Ekaterina Lobacheva, Dmitry Vetrov. Bayesian Compression for Natural Language Processing. In EMNLP 2018.}
% % \end{thebibliography}
	
% \newpage
% \appendix

% \section{Пример секции аппендикса}

\end{document}
